\section{Conclusion}
\label{sec:conclusion}

We asked whether misalignment spreads through populations of interacting LLMs the
way an infection spreads through hosts, and answered it with a measured
transmission kernel fused into a calibrated network epidemic model. Misalignment
propagation \emph{does} follow epidemiological dynamics: it has a measurable
per-contact transmission rate, a recovery process, a basic reproduction number, a
sharp phase transition at the spectral threshold $R_0=(\beta/\mu)\,
\lambda_{\max}(\mathbf{A})=1$, a topology-dependent vulnerability ranking
(scale-free and hub-and-spoke worst), and critical intervention thresholds.

\para{Key takeaway.} The decisive determinant is not network structure but
per-model \emph{susceptibility}: robust models are epidemiologically immune and
confer herd immunity, so whether misalignment self-extinguishes or goes epidemic
is set first by population robustness and only secondarily by topology and
intervention coverage. A population of ${\sim}70$--$80\%$ robust models
self-extinguishes misalignment regardless of how it is wired together. The
cheapest, highest-leverage safety investment is therefore raising individual-model
robustness, followed by protecting hubs in heterogeneous ecosystems.

\para{Future work.} We see five directions: (1) measure the weight-merge
transmission channel directly with GPU fine-tuning to close the
in-context-to-weights gap; (2) characterize the immune/susceptible split across
many model families to estimate the real immune fraction; (3) test switching and
dynamic federated-learning topologies via the joint-spectral-radius threshold; (4)
broaden the evaluation battery and use multiple judges; and (5) formulate the
optimal control problem of where, when, and how much to re-align under a budget.
